\section{DHA Vampire}

The DHA (de Havilland Australia) Vampire is day-fighter and fighter-bomber and also an advanced trainer. It is derived from the de Havilland Vampire, but notably the fighter and fighter-bomber versions use the more powerful Nene engine.

\subsection{Versions}

\subsubsection{DHA Vampire F.30}

The DHA Vampire F.30 is derived from the Vampire F.2, which had trialed the Nene engine in the F.1 airframe. It was equipped with a Nene-2VH engine with 5,000 lb of thrust, significantly more than the Goblin 3 with 3,500 lb. This greater thrust required greater airflow than could be provided by the standard wing-root intakes, so the F.30 was delivered with additional intakes on the upper side of the fuselage behind the cockpit. To improve handling close to the critical Mach number, these were later moved to the lower side of the fuselage.

The F.30 served in the RAAF from 1949 until 1960, when it was replaced by the CAC Sabre. 

\subsubsection{DHA Vampire FB.31}

The DHA Vampire FB.31 is derived from the F.30 with the air-to-ground improvements of the FB.5.

The FB.31 served in the RAAF from  1952 until 1960, when it was replaced by the CAC Sabre.

\subsubsection{DHA Vampire T.33 and T.33A}

The Vampire T.33 is a trainer largely similar to the Vampire T.11 (and notably using the Goblin 35 engine rather than the Nene). The ejection seats and canopy of the T.35 were refitted to the T.33 to give the T.33A.

The T.33 served in the RAAF from 1952 until 1970, when it was replaced by the Aermacchi MB-326H (CAC CA-30). The T.33A conversions took place some time after 1957.

\subsubsection{DHA Vampire T.34 and T.34A}

The T.34 and T.34A are similar to the T.33 and T.33A, but for the RAN rather than the RAAF.

The T.34 served in the RAN from 1954 until 1970, when it was replaced by the Aermacchi MB-326H (CAC CA-30). The T.34A conversions took place some time after 1957.

\subsubsection{DHA Vampire T.35}

The T.35 is a development of the T.33 and adds ejection seats, a new canopy, and increase fuel capacity. The ejection seats and canopy were refitted to the T.33 and T.34 to give the T.33A and T.34A versions.

The T.35 served in the RAAF from 1957 until 1970, when it was replaced by the Aermacchi MB-326H (CAC CA-30).

\subsection{Armament and Stores}

All versions have a gun armament of four Hispano 20 mm cannons with 150 rounds per gun. 

A typical air-to-air load would be two 450L (100 imperial gallon) fuel tanks to increase endurance. 

A typical air-to-ground load would be eight RP-3 rockets and then either two 500 lb bombs or fuel tanks, depending on the mission radius. On short-range missions, two 1,000-lb bombs could be carried but without rockets.

\subsection{ADCs}

\begin{adclist}
    \adcitem{DHA Vampire F.30}
    \adcitem{DHA Vampire FB.31}
    \adcitem{DHA Vampire T.33}
    \adcitem{DHA Vampire T.33A}
    \adcitem{DHA Vampire T.34}
    \adcitem{DHA Vampire T.34A}
    \adcitem{DHA Vampire T.35}
\end{adclist}

\subsection{See Also}

\begin{typelist}
    \typeitem{de Havilland Vampire}
    \typeitem{SNCASE Mistral}
\end{typelist}